\documentclass[10pt,a4paper]{amsart}
\usepackage{amsmath,amssymb,amsthm,graphicx}
% Retained minimal preamble for the Stacks Project chapter more-algebra.tex (More on Algebra).
% The chapter file pinned-original-source.tex carries only \input{preamble} on its first line, so no
% preamble text lives inside the chapter itself.  The Stacks Project uses one single shared preamble for
% every chapter; the snapshot stages it identically as sources/stacks-chow/preamble.tex and
% sources/stacks-etale-cohomology/preamble.tex.  The block below is copied mechanically, line for line,
% from sources/stacks-chow/preamble.tex lines 14-18 (the Xy-pic setup) and 154-217 (the native theorem
% declarations together with the mathematical macros).  Nothing was added, removed or redefined; the
% wrapper class and the amsmath/amssymb/amsthm/graphicx base packages are supplied by the assembler.
% Chapter: more-algebra.tex -- More on Algebra.
% For commutative diagrams we use Xy-pic
\usepackage[all]{xy}

% We use 2cell for 2-commutative diagrams.
\xyoption{2cell}

% Theorem environments.
%
\theoremstyle{plain}
\newtheorem{theorem}[subsection]{Theorem}
\newtheorem{proposition}[subsection]{Proposition}
\newtheorem{lemma}[subsection]{Lemma}

\theoremstyle{definition}
\newtheorem{definition}[subsection]{Definition}
\newtheorem{example}[subsection]{Example}
\newtheorem{exercise}[subsection]{Exercise}
\newtheorem{situation}[subsection]{Situation}

\theoremstyle{remark}
\newtheorem{remark}[subsection]{Remark}
\newtheorem{remarks}[subsection]{Remarks}

\numberwithin{equation}{subsection}

% Macros
%
\def\lim{\mathop{\mathrm{lim}}\nolimits}
\def\colim{\mathop{\mathrm{colim}}\nolimits}
\def\Spec{\mathop{\mathrm{Spec}}}
\def\Hom{\mathop{\mathrm{Hom}}\nolimits}
\def\Ext{\mathop{\mathrm{Ext}}\nolimits}
\def\SheafHom{\mathop{\mathcal{H}\!\mathit{om}}\nolimits}
\def\SheafExt{\mathop{\mathcal{E}\!\mathit{xt}}\nolimits}
\def\Sch{\mathit{Sch}}
\def\Mor{\mathop{\mathrm{Mor}}\nolimits}
\def\Ob{\mathop{\mathrm{Ob}}\nolimits}
\def\Sh{\mathop{\mathit{Sh}}\nolimits}
\def\NL{\mathop{N\!L}\nolimits}
\def\CH{\mathop{\mathrm{CH}}\nolimits}
\def\proetale{{pro\text{-}\acute{e}tale}}
\def\etale{{\acute{e}tale}}
\def\QCoh{\mathit{QCoh}}
\def\Ker{\mathop{\mathrm{Ker}}}
\def\Im{\mathop{\mathrm{Im}}}
\def\Coker{\mathop{\mathrm{Coker}}}
\def\Coim{\mathop{\mathrm{Coim}}}

% Boxtimes
%
\DeclareMathSymbol{\boxtimes}{\mathbin}{AMSa}{"02}

%
% Macros for moduli stacks/spaces
%
\def\QCohstack{\mathcal{QC}\!\mathit{oh}}
\def\Cohstack{\mathcal{C}\!\mathit{oh}}
\def\Spacesstack{\mathcal{S}\!\mathit{paces}}
\def\Quotfunctor{\mathrm{Quot}}
\def\Hilbfunctor{\mathrm{Hilb}}
\def\Curvesstack{\mathcal{C}\!\mathit{urves}}
\def\Polarizedstack{\mathcal{P}\!\mathit{olarized}}
\def\Complexesstack{\mathcal{C}\!\mathit{omplexes}}
% \Pic is the operator that assigns to X its picard group, usage \Pic(X)
% \Picardstack_{X/B} denotes the Picard stack of X over B
% \Picardfunctor_{X/B} denotes the Picard functor of X over B
\def\Pic{\mathop{\mathrm{Pic}}\nolimits}
\def\Picardstack{\mathcal{P}\!\mathit{ic}}
\def\Picardfunctor{\mathrm{Pic}}
\def\Deformationcategory{\mathcal{D}\!\mathit{ef}}

% --- appended 2026-10-12 (Chief): the project own definitions of the reference/slogan/history
% display environments, copied mechanically from the shared Stacks preamble lines 8-12
% (\usepackage{verbatim} plus the three \newenvironment lines, which the chapter preamble keeps
% verbatim from the project preamble). Nothing added, removed or redefined: these environments
% carry no mathematical content and the project itself typesets them as comments. Without these
% three lines a counted statement whose own text uses them fails with Environment slogan undefined.
\usepackage{verbatim}
\newenvironment{reference}{\comment}{\endcomment}
\newenvironment{slogan}{\comment}{\endcomment}
\newenvironment{history}{\comment}{\endcomment}
\usepackage{cleveref}
\crefname{originallemma}{original lemma}{original lemmas}
\Crefname{originallemma}{Original lemma}{Original lemmas}
\makeatletter
\newlabel{algebra-lemma-etale-over-field}{{00U3}{source}{}{sourcepointer-1}{}}
\newlabel{algebra-lemma-etale-over-field@cref}{{[originallemma][0][]00U3}{[1][0][]source}{}{}{}}
\makeatother
\crefname{originallemma}{original lemma}{original lemmas}
\Crefname{originallemma}{Original lemma}{Original lemmas}
\makeatletter
\newlabel{lemma-absolutely-flat-over-absolutely-flat}{{30323}{source}{}{sourcepointer-2}{}}
\newlabel{lemma-absolutely-flat-over-absolutely-flat@cref}{{[originallemma][0][]30323}{[1][0][]source}{}{}{}}
\makeatother
\crefname{originallemma}{original lemma}{original lemmas}
\Crefname{originallemma}{Original lemma}{Original lemmas}
\makeatletter
\newlabel{lemma-absolutely-flat-over-field}{{30504}{source}{}{sourcepointer-3}{}}
\newlabel{lemma-absolutely-flat-over-field@cref}{{[originallemma][0][]30504}{[1][0][]source}{}{}{}}
\makeatother
\crefname{originallemma}{original lemma}{original lemmas}
\Crefname{originallemma}{Original lemma}{Original lemmas}
\makeatletter
\newlabel{lemma-base-change-weakly-etale}{{30303}{source}{}{sourcepointer-4}{}}
\newlabel{lemma-base-change-weakly-etale@cref}{{[originallemma][0][]30303}{[1][0][]source}{}{}{}}
\makeatother
\crefname{originallemma}{original lemma}{original lemmas}
\Crefname{originallemma}{Original lemma}{Original lemmas}
\makeatletter
\newlabel{lemma-composition-weakly-etale}{{30347}{source}{}{sourcepointer-5}{}}
\newlabel{lemma-composition-weakly-etale@cref}{{[originallemma][0][]30347}{[1][0][]source}{}{}{}}
\makeatother
\title{A weakly etale algebra over a field is reduced, integral, a filtered colimit of finite separable products, and closed under subalgebras, quotients and tensor products}
\author{The Stacks Project Authors}
\date{stacks-project:more-algebra.tex:local-snapshot-20261012; GNU FDL 1.2 or later}
\begin{document}
\maketitle
\setcounter{subsection}{1}
\noindent\textit{Editorial note.} All delivered mathematical text is the author's own, copied from the pinned Stacks Project chapter file more-algebra.tex; no statement or proof step was rewritten, reordered or removed.\par\medskip
\noindent\textit{Editorial note.} The chapter file pinned-original-source.tex carries only \textbackslash{}input\{preamble\} on its first line, so no preamble text lives inside the chapter. The retained file sources/stacks-more-algebra/retained-preamble.tex is a mechanical, line-for-line copy of the project's single shared preamble as staged in the snapshot; nothing was added, removed or redefined and no package was added.\par\medskip
\noindent\textit{Editorial note.} The source URL is the source-level Stacks Project URL, as required when more than one card is produced from one source; the per-claim tag is carried by claim\_key.\par\medskip
\noindent\textit{Editorial note.} The author numbers the parts as displayed, and in the proof the sentence 'Part (3) implies (2).' follows the discussion of part (3); that ordering is the author's and is reproduced. The author's phrasing 'any sub or quotient K-algebra' appears in the statement, reproduced unchanged. The six references inside the counted proof span are recorded as non-load-bearing navigation pointers.\par\medskip
\noindent\textit{Allowed standard prerequisites (editorial).}\begin{itemize}
\item Characterisation of weakly etale algebras over a field as reduced algebras integral over the field whose finitely generated subalgebras are etale. invoked at proof lines 30864-30867 and 30873-30877. Is the criterion that turns the finitely generated description into weak etaleness for subalgebras.
\item Etale K-algebras as finite products of finite separable extensions of K. invoked at proof lines 30865-30867. Supplies the shape of the finitely generated subalgebras and of the quotients whose colimits are taken.
\item Permanence of weak etaleness under base change and composition. invoked at proof lines 30887-30888. Gives part (6) for tensor products.
\end{itemize}\medskip
\ifdefined\hypertarget\hypertarget{sourcepointer-1}{}\fi
\noindent\textit{Source reference.} Original lemma 00U3: focused-seeds/AGP/sources/stacks-snapshot-20261011/tags.txt, source lines 1097--1097. Cross-chapter navigation pointer into the Stacks Project; the standard fact it supplies is applied inside the counted proof span and the argument resting on it stays there.\par\medskip
\ifdefined\hypertarget\hypertarget{sourcepointer-2}{}\fi
\noindent\textit{Source reference.} Original lemma 30323: sources/stacks-more-algebra/pinned-original-source.tex, source lines 30323--30323. Same-chapter navigation pointer; the statement it names is standard for the section and the argument that consumes it stays inside the counted proof span.\par\medskip
\ifdefined\hypertarget\hypertarget{sourcepointer-3}{}\fi
\noindent\textit{Source reference.} Original lemma 30504: sources/stacks-more-algebra/pinned-original-source.tex, source lines 30504--30504. Same-chapter navigation pointer; the statement it names is standard for the section and the argument that consumes it stays inside the counted proof span.\par\medskip
\ifdefined\hypertarget\hypertarget{sourcepointer-4}{}\fi
\noindent\textit{Source reference.} Original lemma 30303: sources/stacks-more-algebra/pinned-original-source.tex, source lines 30303--30303. Same-chapter navigation pointer; the statement it names is standard for the section and the argument that consumes it stays inside the counted proof span.\par\medskip
\ifdefined\hypertarget\hypertarget{sourcepointer-5}{}\fi
\noindent\textit{Source reference.} Original lemma 30347: sources/stacks-more-algebra/pinned-original-source.tex, source lines 30347--30347. Same-chapter navigation pointer; the statement it names is standard for the section and the argument that consumes it stays inside the counted proof span.\par\medskip
\begin{lemma}
\label{lemma-class-weakly-etale-over-field}
Let $K$ be a field. If $B$ is weakly \'etale over $K$, then
\begin{enumerate}
\item $B$ is reduced,
\item $B$ is integral over $K$,
\item any finitely generated $K$-subalgebra of $B$ is a finite product
of finite separable extensions of $K$,
\item $B$ is a field if and only if $B$ does not have nontrivial idempotents
and in this case it is a separable algebraic extension of $K$,
\item any sub or quotient $K$-algebra of $B$ is weakly \'etale over $K$,
\item if $B'$ is weakly \'etale over $K$, then $B \otimes_K B'$ is
weakly \'etale over $K$.
\end{enumerate}
\end{lemma}

\begin{proof}
Part (1) follows from Lemma \ref{lemma-absolutely-flat-over-absolutely-flat}
but of course it follows from part (3) as well.
Part (3) follows from Lemma \ref{lemma-absolutely-flat-over-field}
and the fact that \'etale $K$-algebras are finite products of
finite separable extensions of $K$, see
Algebra, Lemma \ref{algebra-lemma-etale-over-field}.
Part (3) implies (2).
Part (4) follows from (3) as a product of fields is
a field if and only if it has no nontrivial idempotents.

\medskip\noindent
If $S \subset B$ is a subalgebra, then it is the filtered
colimit of its finitely generated subalgebras which are
all \'etale over $K$ by the above and hence
$S$ is weakly \'etale over $K$ by
Lemma \ref{lemma-absolutely-flat-over-field}.
If $B \to Q$ is a quotient algebra, then
$Q$ is the filtered colimit of $K$-algebra quotients of
finite products $\prod_{i \in I} L_i$ of finite separable extensions
$L_i/K$. Such a quotient is of the form $\prod_{i \in J} L_i$
for some subset $J \subset I$ and hence the result
holds for quotients by the same reasoning.

\medskip\noindent
The statement on tensor products follows in a similar manner
or by combining Lemmas \ref{lemma-base-change-weakly-etale} and
\ref{lemma-composition-weakly-etale}.
\end{proof}

\end{document}
